\chapter{R\'ealisation}
\label{chap:realisation}

\section{Introduction}
% Annoncer le contenu : environnement de travail, technologies,
% fonctionnalites realisees (avec captures d'ecran), tests.

[\`A compl\'eter]

\section{Environnement de travail}

\subsection{Environnement mat\'eriel}
[\`A compl\'eter : caract\'eristiques de la machine de d\'eveloppement.]

\subsection{Environnement logiciel}
% Lister les outils utilises : IDE (VS Code), Node.js, Angular CLI,
% gestionnaire de version (Git/GitHub), Postman, navigateur, etc.

\begin{itemize}
    \item \textbf{Visual Studio Code} : \'editeur de code ;
    \item \textbf{Node.js / npm} : environnement d'ex\'ecution et gestion
    des d\'ependances ;
    \item \textbf{Angular CLI 21} : g\'en\'eration et build du projet ;
    \item \textbf{Git / GitHub} : gestion de versions ;
    \item [\`A compl\'eter : Postman, Docker, IDE backend, etc.]
\end{itemize}

\section{Technologies et biblioth\`eques utilis\'ees}
% Justifier les choix technologiques principaux (cf. package.json et
% DOCUMENTATION.md section 2).

\begin{table}[H]
    \centering
    \begin{tabularx}{\textwidth}{l l X}
        \toprule
        \textbf{Technologie} & \textbf{Version} & \textbf{R\^ole} \\
        \midrule
        Angular & 21.2 & Framework frontend (composants standalone) \\
        TypeScript & 5.9 & Langage principal du frontend \\
        Keycloak JS & 26.2 & Authentification SSO et gestion des r\^oles \\
        bpmn-js & 18.16 & Visualisation des diagrammes BPMN 2.0 \\
        RxJS & 7.8 & Programmation r\'eactive (appels HTTP, flux de donn\'ees) \\
        jsPDF / jspdf-autotable & -- & G\'en\'eration de rapports PDF \\
        xlsx (SheetJS) & 0.18 & Export des donn\'ees au format Excel \\
        \bottomrule
    \end{tabularx}
    \caption{Principales technologies et biblioth\`eques utilis\'ees (frontend)}
    \label{tab:technologies}
\end{table}

[\`A compl\'eter : technologies c\^ot\'e backend (Spring Boot, Drools, base
de donn\'ees, etc.) et justification des choix.]

\section{Travail r\'ealis\'e}
% Pour chaque module, presenter brievement son fonctionnement et
% illustrer avec 1-3 captures d'ecran commentees.

\subsection{Authentification et gestion des utilisateurs}
[\`A compl\'eter : description + captures d'\'ecran (page de connexion,
inscription, gestion des utilisateurs).]

\subsection{Gestion des processus m\'etier (BPMN)}
[\`A compl\'eter : description + captures d'\'ecran (liste des processus,
gestion des t\^aches, visualiseur BPMN, export du fichier .bpmn).]

\subsection{Moteur de r\`egles m\'etier}
[\`A compl\'eter : description + captures d'\'ecran (liste des r\`egles,
formulaire de cr\'eation, aper\c{c}u DRL, historique des versions).]

\subsection{Module IA / NLP}
% S'appuyer sur DOCUMENTATION_IA_NLP.md (pipeline NLP, exemples de
% phrases -> regles -> DRL, score de confiance, gestion des
% ambiguites) pour illustrer ce module en detail : c'est probablement
% la partie la plus valorisable du rapport.

[\`A compl\'eter : description du pipeline NLP (tokenisation,
reconnaissance d'entit\'es, g\'en\'eration DRL, score de confiance) +
captures d'\'ecran des modes Saisie / Aper\c{c}u / \'Edition + exemples
concrets de phrases converties.]

\subsection{Import / Export de dossiers}
[\`A compl\'eter : description + captures d'\'ecran (formulaires de
saisie, listes filtr\'ees, exports PDF/Excel).]

\section{Tests}
% Decrire les tests realises : tests unitaires (Vitest), tests
% manuels/fonctionnels par cas d'utilisation, eventuels tests
% d'integration avec le backend.

[\`A compl\'eter : strat\'egie et exemples de tests (unitaires avec Vitest,
tests fonctionnels manuels par sc\'enario).]

\section{Conclusion}
% Resumer le travail realise et faire la transition vers la conclusion
% generale.

[\`A compl\'eter]

\clearpage
